\section{Proof for the end attachments}\label{sec:attachments}

We prove Theorem~\ref{attach:stable}, using the varieties and root lists
defined in Section~\ref{sec:bundle-families}.

\subsection{Geometry and divisor degrees}
The lattice is the direct sum of the hexagon lattices and the base
lattices. In a base factor of dimension $d$, the lifted rays are
$e_1,\ldots,e_d,t-\sum e_i$, where $t$ is its stated twisting vector.
Maximal cones contain adjacent ray pairs in each hexagon and omit one
ray in each base factor. The simplex lengths in the anticanonical polytope are at least one,
so its normal fan is precisely the displayed complete bundle fan;
the block triangular maximal-cone matrices are unimodular. The primitive
collections are the hexagon primitive collections and the lifted
primitive collections of the base. Each leaf relation has degree
$2-1=1$, and each internal base relation has degree $3-2=1$.
The fibre primitive relations have degrees one and two. Thus the
variety is Fano. The incidence graph is connected; within each hexagon
the rays form one matroid component, and each lifted base relation
connects its base coordinates to all the incident hexagon components.
Hence the ray matroid is connected.

The number of base factors is $j+a+b+1$ and their total dimension is
$2j+a+b$. The dimension and Picard formulas follow. We prove stability
using the following explicit degree and rank calculations.
\medskip
The dual anticanonical hexagon is
\[
 Q=\{(x,z): |x|,|z|,|x+z|\le1\}.
\]
Its vertices, cyclically ordered, are
$(1,0),(1,-1),(0,-1),(-1,0),(-1,1),(0,1)$.
For $A_a=(t_1,\ldots,t_{a+1})$ set
\[
 P_a(x,z)=\prod_{\ell=1}^{a+1}(2+\langle(x,z),t_\ell\rangle),
 \qquad
 m_a(p)=\int_Q x^pP_a(x,z)\,dx\,dz\quad(0\le p\le2).
\]
The three-entry moment vectors are
\[
\begin{array}{c|ccc}
 a&m_a(0)&m_a(1)&m_a(2)\\\hline
 0&6&-5/6&5/3\\
 1&139/12&-5/3&10/3\\
 2&67/3&-31/10&31/5\\
 3&43&-86/15&1373/120
\end{array}
\]
For a hexagon ray $h$, define $m_{a,h}$ by the same formula integrated
over its facet $\langle(x,z),h\rangle=-1$ with lattice measure.
For either ray of leaf $\ell$, define $m_{a,\ell}$ by omitting the
factor $2+\langle(x,z),t_\ell\rangle$ from $P_a$.
These prescriptions specify every modified moment vector without
rounding. Reflection changes a right-hand moment vector $v_p$ to
$(-1)^pv_p$. The reflected local hexagon ray $h$ is the global ray $-h$.

Write
\[
 W=\begin{pmatrix}3&0&5/6\\0&5/6&0\\5/6&0&7/15\end{pmatrix},\quad
 \kappa=\begin{pmatrix}9/2&-3&1/2\\3&-1&0\\1/2&0&0\end{pmatrix},\quad
 \kappa_1=\begin{pmatrix}3&-1&0\\1&0&0\\0&0&0\end{pmatrix}.
\]
These are respectively the moments of $(2-|x|)dx$, the coefficients
of $(3+x-x')^2/2$, and the coefficients of $3+x-x'$.
For $j\ge2$, let $L_0=m_a$ be a row vector and let $R_0$ be the
reflected $m_b$ as a column vector. Set
\[
 L_{i+1}=L_i\kappa W,\qquad R_{i+1}=W\kappa R_i.
\]
The anticanonical polytope volume is
$Z=L_0\kappa R_{j-2}$, and $(-K_X)^n=n!Z$.
Normalized weights are facet volumes divided by $Z$.
For the left end replace $L_0$ by its modified vector; for the right
replace $R_0$ by its reflected modified vector. For the internal base
factor $c$, $1\le c\le j-1$, each of its three rays has weight
\[
 \beta_c=\frac{L_{c-1}\kappa_1R_{j-1-c}}{Z}.
\]
For an internal hexagon $i$, $1\le i\le j-2$, replace $W$ in
$L_{i-1}\kappa W\kappa R_{j-2-i}$ by its facet moment matrix $W_h$.
Its denominator is $Z$. These formulas follow by integrating the
simplex slices and applying Fubini. The case $j=1$ is simply the
integral of $P_a(x,z)P_b(-x,-z)$ over $Q$.

\subsection{End subspaces and the rank recurrence}
An end of type $a$ has lattice $E_a=\ZZ^2\oplus\ZZ^{a+1}$ and rays
\[
 (h,0)\ (h\in\{\pm r,\pm q,\pm(r+q)\}),\quad
 (0,e_\ell),\ (t_\ell,-e_\ell)\quad(1\le\ell\le a+1).
\]
Enumerate the ray-spanned subspaces by starting with zero, adjoining
each ray in turn, and closing under inclusion of all rays in the span.
This finite, exact linear-algebra procedure gives respectively
$13,34,92,241$ subspaces for $a=0,1,2,3$, including zero and the whole
space. All their ray subsets and ranks are supplied with the computation.
For such an end subspace $S$, record its rank, selected ray weight,
whether $r\in S$, and whether $S+\CC r=E_{a,\CC}$.
No symmetry of either end is assumed.

Use only the independent reflections of the internal hexagons fixing
$r_i$, and the transpositions of the two untwisted rays of each internal
$\PP^2$. These preserve the entire fan, including the end pieces.
Every invariant ray-spanned subspace is represented by an end choice
on each side, choices $H_i\in\{0,1,2\}$ of empty set, distinguished line,
or whole plane in each internal hexagon, and binary choices $C_c,T_c$
for the two untwisted base rays and the twisted ray. Passing to closure
can only increase weight at fixed rank, so minimizing over all these
configurations suffices.

The outgoing boundary state is $G$ if its distinguished ray was selected,
$g$ if that line is spanned through the preceding relations, and $u$ if
it is not spanned. Initialize at the left end with $G$ when $r\in S$
and $u$ otherwise, with rank $\dim S$. Adjoining an internal base and
the following hexagon adds the rank
\[
 d=2C+T(1-C)+
 \begin{cases}
 H+\mathbf1_{CT=1,\,s=u},&H\ge1,\\
 CT,&H=0,
 \end{cases}
\]
and gives state $G$ when $H\ge1$, state $g$ when $H=0$, $CT=1$,
and $s\ne u$, and state $u$ otherwise. Its selected weight is
$(2C+T)\beta+\omega(H)$, where $\omega(H)$ is the empty, line,
or plane weight. This is the usual path relation
$-g_c+r_{c-1}-r_c$ modulo previously selected subspaces.

For the selected subspace $V$ in the processed ambient space $E$, define
binary indicators $\eta=\mathbf1_{V\ne0}$ and
$F=\mathbf1_{E/(V+\CC r)=0}$, where $r$ is the outgoing distinguished
ray. Initially $F=\mathbf1_{S+\CC r=E_{a,\CC}}$. At an internal step,
$\eta'=\mathbf1_{\eta=1\text{ or }d>0}$ and
\[
 F'=F\ \text{and}\ \bigl(\mathbf1_{s=u}+4-d
                              =\mathbf1_{s'=u}\bigr).
\]
Indeed, if $F$ holds, the old rank deficiency is exactly
$\mathbf1_{s=u}$. Otherwise, a missing direction independent of the
boundary can never be supplied by the remaining chain. This also proves
that $F=0$ implies $F'=0$.

To close with right end subspace $S_R$, let $e=\dim S_R$ and let
$g_R$ mean $r\in S_R$. The last base and right end add rank
\[
 d_R=2C+T(1-C)+e+\mathbf1_{CT=1}\mathbf1_{s=u\ \mathrm{or}\ \neg g_R}.
\]
The whole configuration has full rank precisely when
\[
 F\ \text{and}\ \bigl(\mathbf1_{s=u}+2+(b+3)-d_R=0\bigr).
\]
The formula for $d_R$ follows because the remaining relation
$r_{j-2}-r_{j-1}$ is dependent exactly when both boundary lines are
already spanned. Minimize rank minus selected weight, retaining only
the least cost for each state and pair of binary indicators. Discard the empty
and full-rank final configurations. This proves the finite recurrence;
it is not inferred from comparisons with numerical ranks.

\subsection{Uniform comparison and stability}
The internal kernel is positive and has Hilbert projective contraction
factor $\tau=2/7$ and diameter $\Delta=\log(81/25)$, as in the path theorem.
The end polynomials are positive on $Q$. Equivalently, their moment
vectors are those of the positive density
\[
 g_a(x)=\frac{1}{2-|x|}\int_{\{z:(x,z)\in Q\}}P_a(x,z)\,dz
\]
against $(2-|x|)dx$. After one application of the kernel the outgoing
functions have oscillation ratio at most nine.
Thus they converge to the same limiting functions as the original path
ends. Every internal hexagon weight is at most $27$, and every base
weight is at most two. At an end, $1\le P_a\le3^4=81$ and the incoming
function has oscillation ratio at most nine. Since a hexagon facet has
lattice length one and $\operatorname{vol}(Q)=3$, every end hexagon
weight is at most $81\cdot9/3=243$. End leaf weights are at most one.
These bounds also hold for limiting and finite comparison weights.

Take $q_0=20$, comparison index $K=40$, and $j=2q_0+m$.
The left end and periods $1,\ldots,q_0-1$ retain their left chain
indices and use the limiting right function. Periods
$q_0,\ldots,q_0+m-1$ use both limiting functions. The remaining
periods and the right end retain their right indices and use the
limiting left function. There are at most $9q_0+10$ rays in either
end zone. Write
\[
 \bar\Delta=11757/10000,\quad x=\bar\Delta\tau^{q_0-3},\quad
 y=2\bar\Delta\tau^{K-2},\quad B=243.
\]
The deliberately weaker exponents cover the first kernel application
from an end piece, including its adjacent base factor. The contraction
bound and $e^t-1\le t+t^2$ imply the uniform true-path error
\[
 E_{\rm path}=B\left\{9\frac{2x}{1-\tau}(1+2x)
                    +2(9q_0+10)(x+x^2)\right\}.
\]
For each of the four comparison paths $m=0,1,2,3$, replacing infinite
indices by $K$ changes every configuration cost by at most
\[
 E_{\rm rat}=B\{9(2q_0+4)+10\}(y+y^2).
\]
This also covers every simple internal transition cycle.

The limiting all-selected period has weight four: summable convergence
errors make the difference between total weight and the number of
periods times the limiting weight bounded independently of $j$, whereas
the total weight is $4j+a+b$. Hence the all-selected $G$ loop and empty
$u$ loop have exactly zero cost. All other simple cycles are the same
as in the original path theorem and have cost exceeding $0.1259$.
The exact rational comparison calculation gives the following strict
lower bounds; these are terminating decimal lower bounds, not rounded
minima:
\begin{table}[ht]
\centering
\caption{Strict lower bounds for nonempty proper-rank configurations.}
\label{attach:bounds}
\begin{adjustbox}{max width=\textwidth,center}
\begin{tabular}{ccc}
\toprule
$(a,b)$ & every true $2\le j\le44$ & each comparison $m=0,1,2,3$\\
\midrule
\csname @@input\endcsname tables/attachment_bounds.tex
\bottomrule
\end{tabular}
\end{adjustbox}
\end{table}
Their all-selected comparison costs have absolute value less than
$6\cdot10^{-25}$. The exact rational bounds give
\[
 E_{\rm path}+2E_{\rm rat}+6\cdot10^{-25}<0.00006523.
\]
The finite inputs are fully specified by the displayed rays, moment
integrals, matrices, recurrence, and constants. The supplied code uses
rational arithmetic to evaluate them. The $j=1$ cases are checked by
positive rational decompositions into ray bases with strongly connected
exchange graphs, as in Lemma~\ref{lem:basecertificate}. Their normalized
degrees independently agree with direct integration of the product of
the stated linear factors over $Q$ and its facets. The decompositions
are supplied in the ancillary data.

\begin{proof}[Completion of the proof of Theorem~\ref{attach:stable}]
Apply the cycle-deletion argument of Appendix~\ref{sec:paths} to the
central periods. A walk between repeated states either contains a positive-cost
simple cycle, or consists solely of one of the two zero loops. Deleting
all-selected $G$ loops preserves nonemptiness and proper rank. Deleting
empty $u$ loops does likewise: the boundary at which deletion occurs is
an internal hexagon, so state $u$ implies that this hexagon's independent
nondistinguished direction is missing; the shortened configuration
cannot become full rank. Thus the end pieces cause no new exception in
the zero-loop argument. A full-rank reduced configuration has cost at
least that of the all-selected configuration, by positivity of weights.
Deleting cycles therefore reduces every comparison configuration to
$m\le3$, or supplies the positive cycle bound. Subtracting all errors,
every nonempty proper-rank configuration at arbitrary length has deficit
\[
 >0.028404-0.00006523=0.02833877>0.
\]
Together with the finite cases and connectedness of the ray matroid,
the symmetry criterion proves the theorem.
\end{proof}
